\begin{prob}
    In each of the problems below state the best case and worst case time
    complexities of the given piece of code using asymptotic notation. Note that
    some algorithms may have the same best case and worst case time complexities. If
    the best and worst case complexities \emph{are} different, identify which
    inputs result in the best case and worst case. You do not otherwise need to show your work for
    this problem.

    \emph{Example Algorithm:} \python{linear_search} as given in lecture. 

    \emph{Example Solution:}
    Best case: $\Theta(1)$, when the target is the first element of the array.
    Worst case: $\Theta(n)$, when the target is not in the array.

    \begin{subprobset}

        \begin{subprob}
            \begin{minted}[autogobble]{python}
                def kth_largest(numbers, k):
                    """Finds the k-th largest element in the array.
                    `numbers` is an array of n numbers."""
                    n = len(numbers)
                    for i in range(n):
                        count = 0 # Count how many numbers are larger than numbers[i]
                        for j in range(n):
                            if numbers[j] > numbers[i]:
                                count += 1
                        if count == k - 1:
                            return numbers[i]
            \end{minted}
            
            \begin{soln}
                Best case: $\Theta(n)$, when the first element of the array is the $k$th largest.

                Worst case: $\Theta(n^2)$, we need to count the number of elements larger than each element. If the last element is $k$th largest, then we need to run entire nested loops.
            \end{soln}
        \end{subprob}

        \begin{subprob}
            \begin{minted}[autogobble]{python}
                def index_of_kth_largest(numbers, k):
                    """`numbers` is an array of size n"""
                    # the kth_largest() from above
                    element = kth_largest(numbers, k)
                    # the linear_search() from lecture 6 slide 11
                    return linear_search(numbers, element)
            \end{minted}
            

            \begin{soln}
                Best case: $\Theta(n)$. This occurs when the $k$th largest element is the first element of the array,
                as then \python{kth_largest} takes $\Theta(n)$ and \python{linear_search} takes $\Theta(1)$,
                for a total of $\Theta(n)$.

                Worst case: $\Theta(n^2)$. This occurs when the $k$th largest element is the last element of the
                array. In this situation, \python{kth_largest} takes $\Theta(n^2)$ time and \python{linear_search}
                takes $\Theta(n)$ time, for a total of $\Theta(n^2)$ time.
            \end{soln}
        \end{subprob}

    \end{subprobset}


\end{prob}
